% gerado por ferramentas/md2tex_enunciado.py a partir de questoes/q113.md (não editar à mão)
O comprimento da parte oscilante da corda de um violão determina a nota musical emitida. No braço do violão, são colocadas peças metálicas espaçadas chamadas trastes. Quando as cordas são pressionadas entre os trastes, os comprimentos de suas partes oscilantes são reduzidos, gerando ondas de diferentes frequências. Algumas notas com os valores aproximados de suas frequências são apresentadas no quadro. A figura ilustra um violão com cordas de comprimento total vibrante ($L_T$) igual a 60\,cm. Considere o modo de vibração fundamental ($n = 1$) e que as frequências de vibração são inversamente proporcionais ao comprimento da parte oscilante da corda. Sabe-se que a primeira corda oscila com frequência de 660\,Hz (a nota mi) quando não pressionada (corda livre).\par
\begin{figlado}{0.4}{q113_fig1.png}
\textbf{Frequência das notas musicais, em hertz, no sistema temperado}\par
\begin{center}\renewcommand{\arraystretch}{1.15}
\begin{tabular}{cccccccc}
\toprule
\textbf{Oitava} & \textbf{Dó} & \textbf{Ré} & \textbf{Mi} & \textbf{Fá} & \textbf{Sol} & \textbf{Lá} & \textbf{Si}\\
\midrule
3ª & 130 & 147 & 165 & 175 & 196 & 220 & 248\\
4ª & 260 & 294 & 330 & 349 & 392 & 440 & 495\\
5ª & 520 & 587 & 660 & 698 & 784 & 880 & 990\\
6ª & 1\,040 & 1\,174 & 1\,320 & 1\,397 & 1\,568 & 1\,760 & 1\,980\\
\bottomrule
\end{tabular}\end{center}
\end{figlado}
\fonte{OLIVEIRA, Leonardo M.; VEIT, Eliane A.; SCHNEIDER, Claúdio. \textbf{Oitavas e notas}. Disponível em: www.if.ufrgs.br. Acesso em: 6 nov. 2023 (adaptado).}
Ao pressionar a primeira corda, de forma que a sua parte oscilante tenha 45 centímetros, e fazê-la vibrar, qual é a nota emitida pela corda?\par
\begin{altlinha}
\alth{A}{Dó.}%
\alth{B}{Mi.}%
\alth{C}{Sol.}%
\alth{D}{Lá.}%
\alth{E}{Si.}
\end{altlinha}
